

\section{Theoretical background}
Our model exists as a combination of two structures: signalling games and social choice functions. We first describe the latter, and then 


\subsection{Signalling games}
Signalling games are dynamic Bayesian games, first studied in relation to learning by Lewis \citeme. The players have incomplete information about the initial situation, and the game unfolds sequentially over time, with actions known to the other players \citeme. 

More precicely, one player -- often called the sender -- has some private information about the state of the world, an takes an observable action $S$, called the signal, after which another player -- often called the reciever -- updates their beliefs about the world and chooses an action. In our case of deterrence, incomplete information is essential; the uncertainties around the opponents resolve lies at the heart of the mechanism making deterrence a viable strategy. 

\subsection{Social choice functions}






\section{The model}

Our model adopts this basic structure of a signalling game, as presented above, where we assume one of the players to be a coalition of multiple states $\C = \{1, 2, \dots, n\}$. The other player is some potential attacking adversarial actor $\A$. The adversary $\A$ has a \emph{type}, $T \in \{0,1\}$, where $1$ represents being agressive -- actively planning an attack -- and $0$ represents a benign adversary. The coalition $\C$ holds a common prior $\pi = \Pr(T=1) \in (0,1)$, describing their belief that $\A$ has type $1$.

If the adversary $\A$ attacks, i.e. $T=1$, each member $i$ of the coalition $\C$, recieves a binary warning signal $s_i \in \{0,1\}$, giving a probability
\[p_i = \Pr(s_i = 1 \mid T=1).\]
If the adversary is of type $T=0$, then coalition members recieve false alarms with probability
\[q_i = \Pr(s_i = 1 \mid T=0).\]
We assume that signals are independent across members. 

Each coalition member $i\in \C$ then decides on a binary vote $v_i \in \{0,1\}$, interpreted as their vote towards authorizing retaliation. 

In order to focus this paper on the institutional mechanisms, and not on wether the coalition members try to strategize for their own gain, we assume that $v_i = s_i$, which means that we assume the actors to vote truthfully based on their own intelligence assessments. 

Ulike a classical signalling game, the coalition does not decide on a signal. The signal is \emph{implied} by the institutional rule that determines the outcome of the voting process based on the individual coalition members' votes. 

\subsection{Weighted threshold functions}

The coalition is assumed to have a publically known decision rule, formalized by a social choice function $f$. In order to be able to analyze the behaviour of our signalling game with respect to such a function, we need to restrict to a class of viable rules. Fist, we assume a binary outcome, $1$ representing coalition retaliation and $0$ representing no retaliation. Second, we assume the choice function $f$ to be a weighted threshold rule. This means that each coalition member $i$ is assigned a weight $w_i \in \R$ dependent on the context of the attack, and that there is a threshold $\tau \in \R$, such that the aggregated vote of the coalition is $1$ if the weighted sum of the individual members votes exceeds the threshold. This weight scheme is meant to include the possibility that the attacked country in the coalition gets a vote that counts more towards the total, or that frontline countries have a larger say, or that a country has far better ISR capabilities etc. 

Formally, then, given a list of votes coalition individual votes $v=(v_1, v_2, \dots, v_n)$, and a weight scheme $w=(w_1, \dots, w_n)$, we have 
\[
f(v;w) = 
\begin{cases}
    1, &\displaystyle\sum_{i=1}^n w_i v_i \geq \tau \\
    0, &\text{else}    
\end{cases}
\]
Most viable voting rules are formalizable in this framework, like the majority rule, consensus rule, dictatorship, qualified majority, veto players, etc. There are also some decision rules that are not possible to model in this framework, like multi-stage rules, lexicographical rules, non-monotonic rules and randomized rules. However, such rules are very unlikely to be viable for choosing a retaliation action, hence our focus on weighted threshold rules. 

\subsection{Retaliation probabilities}

Assuming the attacker type is $T=1$, retaliation will occur if the weighted number of coalition members recieving the signal meets the threshold. Because these signals are independent, the probability of retaliation is 
\[R(\tau) = \Pr\left(\sum_{i=1}^n w_i v_i \geq \tau \mid T=1\right).\]
This function captures the \emph{deterrence credibility} of the coalition $\C$. In other words it is the probability that an attack by the adversary $\A$ triggers a retaliatory response. 

Similarly, if the adversary's type is $T=0$, then retaliation is accidental, with probability 
\[F(\tau) = \Pr\left(\sum_{i=1}^n w_i v_i \geq \tau \mid T=0\right).\]
This function captures the coalition's \emph{escalation risk}, i.e., the probability that the coalition accidentaly authorizes a response when there had been no attack. This could happen if, for example, another actor than $\A$ was responsible for an attack on the coalition, or that an accident is interpreted as an attack. 

Notice that a lowe threshold $\tau$ increases the probability of both deterrence and accidental use. 

In the simple case that each coalition member $i$ has an equal weight, say $w_i = 1$ for all $i$, as well as equal probabilities $p_i$, and the threshold $\tau$ is a positive integer less than or equal to $n$, then these probabilities become binomial, with 
\[R(\tau) = \sum_{k=\tau}^{n} \binom{n}{k} p^{k} (1-p)^{n-k},\]
and similarly 
\[F(\tau) = \sum_{k=\tau}^{n} \binom{n}{k} q^{k} (1-q)^{n-k}.
\]
The often used majority rules and consensus rules are both of this nature, and these closed formula allows us to study these in slightly more detail. 

\subsection{Costs and benefits}

The optimal outcome of the adversay actor $\A$ is to attack and suffer no consequences. The optimal outcome for the coalition $\C$ is that no attack occurs. 

In order for an attack to be of interest, there has to be a percieved benefit for $\A$. In the situation of optimal outcome for $\A$, where it attacks and no retaliation occurs, we denote this benefit by $B$. We assume this value to be a number, for example monetary benefit or some other measure of benefit that $\A$ deems desirable. If retaliation occurs, then $\A$ suffers a cost $C$, which we assume is of comparable type to the benefit $B$. 

The expected payoff for $\A$ not attacking is naturally $0$. The probability of retaliation by $\C$ is given as above by $R(\tau)$, hence the expected payoff for attacking is 
\[(1-R(\tau))B-R(\tau)C.\]
Assuming rationality, the adversary state attacks if it deems the expected payoff to be better for attacking versus not attacking, in other words if the above value is greater than $0$. Rearanging the terms gives that $\A$ attacks if 
\[R(\tau)\leq \frac{B}{B+C}.\]
This fraction is the adversary actor $\A$'s \emph{deterrence threshold}, and we denote it for simplicity by $\bR := \frac{B}{B+C}$. This means we have two actions:
\begin{enumerate}
    \item If $R(\tau)\leq \bR$, then the adversary actor $\A$ attacks;
    \item If $R(\tau)\geq \bR$, then the coalition $\C$'s collective deterrence has succeeded. 
\end{enumerate}











\section{Equilibria and implications}


